% Generated draft: requires mathematical and source review.
\documentclass[11pt]{article}
\usepackage[margin=25mm]{geometry}
\usepackage{fontspec}
\setmainfont{Latin Modern Roman}
\usepackage{amsmath,amssymb,mathtools,unicode-math}
\setmathfont{Latin Modern Math}
\usepackage{enumitem,needspace,xurl,hyperref}
\setlength{\parindent}{0pt}
\setlength{\parskip}{5pt}
\setlength{\emergencystretch}{3em}
\newcommand{\N}{\mathbb{N}}
\newcommand{\Z}{\mathbb{Z}}
\newcommand{\Q}{\mathbb{Q}}
\newcommand{\R}{\mathbb{R}}
\newcommand{\C}{\mathbb{C}}
\newcommand{\F}{\mathbb{F}}
\newcommand{\A}{\mathbb{A}}
\newcommand{\PP}{\mathbb{P}}
\newcommand{\E}{\mathbb{E}}
\newcommand{\eps}{\varepsilon}
\newcommand{\dd}{\,\mathrm{d}}
\begin{document}
\section*{UnsolvedMath ID: 20001952}
\textbf{Source problem code:} catalogue entry
\subsection{Definitions and mathematical statement}
T. Hausel asks for a Martin nonabelian-to-abelian integration theorem for 3-Sasakian quotients, and for Sasakian quotients, suggesting that Kirwan surjectivity only through middle degree might suffice in the former case. The source does not specify actions, regularity, coefficient field, quotient category or orientations. The supplied aid treats regular free reductions and liftable ordinary cohomology classes, not a theorem for all singular or orbifold quotients.

Here is an explicit conditional correspondence framework, with real cohomology. Let a compact connected Lie group \(G\) act on a manifold \(M\), let \(T\subset G\) be a maximal torus with Weyl group \(W\), and let
\[
 \mu_G:M\to(\mathfrak g^*)^{\oplus r},\qquad r\geq1,
\]
be an equivariant multicomponent moment map. Let \(\mu_T\) be its componentwise projection to \(\mathfrak t^*\). Assume the zero levels are regular, the relevant \(G\)- and \(T\)-actions are free, and the quotient manifolds
\[
 X_G=\mu_G^{-1}(0)/G,\quad X_T=\mu_T^{-1}(0)/T,\quad
 Y=\mu_G^{-1}(0)/T
\]
are compact where needed for integration. There are maps
\(i:Y\hookrightarrow X_T\) and \(p:Y\to X_G\) with fibre \(G/T\). A compatible Martin lift of \(a\in H^*(X_G;\R)\) is \(\widetilde a\in H^*(X_T;\R)\) satisfying \(i^*\widetilde a=p^*a\).

Choose positive roots and orient \(\mathfrak g/\mathfrak t\) by its positive-root complex planes. For the associated root lines over \(X_T\), put
\[
 e_+=\prod_{\beta>0}c_1(L_\beta)
 \in H^{\dim G-\dim T}(X_T;\R).
\]
The formula requires compatible correspondence orientations: the normal bundle of \(Y\) in \(X_T\) is the direct sum of \(r\) copies of the associated root bundle, with Euler class \(i^*e_+^r\), and the vertical tangent bundle of \(p\) has Euler class \(i^*e_+\). With zero-set, fibre and quotient orientations chosen compatibly, the aid's lift-dependent formula is
\[
 \int_{X_G}a=\frac1{|W|}\int_{X_T}\widetilde a\,e_+^{r+1}.
\]
Integration selects top-degree components. The degree matches
\(\dim X_T-\dim X_G=(r+1)(\dim G-\dim T)\), and fibre integration of the vertical Euler class is \(\chi(G/T)=|W|\). If one uses different conventional root or symplectic orientations, the corresponding sign must be tracked; no orientation-free sign convention is inferred from the abbreviated source.

For a Kirwan-lift construction, an equivariant class \(b\in H_G^*(M;\R)\) with \(\kappa_G(b)=a\) yields
\(\widetilde a=\kappa_T(\operatorname{res}_T^G b)\). The formula therefore holds on the liftable image without assuming full Kirwan surjectivity. If \(\kappa_G\) is surjective in degrees through \(q\), it guarantees liftability of the subalgebra generated by \(\bigoplus_{j\leq q}H^j(X_G;\R)\). It may lift additional classes too; middle-range surjectivity does not identify its entire image or force the orientation class by Poincaré duality.

In the regular free 3-Sasakian case \(r=3\), the factor is \(e_+^4\); in the Sasakian/contact case \(r=1\), it is \(e_+^2\). The aid invokes Galicki--Salamon low-odd-cohomology vanishing for a smooth compact 3-Sasakian quotient of dimension \(D=4k+3\): odd cohomology through \(\lfloor D/2\rfloor=2k+1\) vanishes. Hence the subalgebra guaranteed by surjectivity through that range is even-degree, and its degree-\(D\) part is zero. It cannot contain the nonzero orientation class. This identifies a limitation of the suggested middle-degree argument, without disproving full surjectivity or other Martin-type results.

J. Munn's Comment 2.5 proposes an analytic polysymplectic extension using 3-Sasakian manifolds as boundaries of hyperkähler spaces and ``compact cohomology,'' ignoring cone points for integration. This remains a proposal here: compactly supported versus relative cohomology, singular-cone treatment, boundary conditions and analytic convergence are not defined, so no unconditional boundary integration theorem is asserted.
\subsection{Short English statement}
Under regular free correspondence hypotheses, compatible cohomology lifts satisfy a Martin formula with factor \(e_+^4\) for 3-Sasakian reduction and \(e_+^2\) for contact reduction. Middle-degree surjectivity only guarantees an algebra of lifts; low odd-cohomology vanishing prevents that algebra alone from containing the odd-dimensional orientation class.
\subsection{Sources}
\begin{itemize}
\item Permanent UnsolvedMath ID \(20001952\), exact public code catalogue entry; local registry \texttt{aim-geometry-notes.json}, \texttt{notes[289]}, canonical tag question.
\item AIM workshop ``Moment maps and surjectivity in various geometries,'' Question 2.4 attributed to T. Hausel and Comment 2.5 attributed to J. Munn: \url{https://aimath.org/WWN/momentmaps/momentmaps.pdf}. No PDF was accessed. The source's hyperkähler encoding defect is normalized in the prose only.
\item Supplied attributions: Martin's integration theorem and Galicki--Salamon low-odd-cohomology vanishing. Full bibliographic details and the exact theorem hypotheses are absent; no original literature entry or year is supplied.
\item Research attempt 1, ``A lift-dependent Martin formula and a sharp middle-degree obstruction,'' labelled partial\_result, open at its run, proved and low novelty confidence; registry partially\_solved. The verification note says a mathematical transcription was verified but does not include that promised transcription in the supplied field.
\end{itemize}
\end{document}
